\documentclass[11pt]{article}
\usepackage[final]{acl}
\usepackage{times}
\usepackage{latexsym}
\usepackage[T1]{fontenc}
\usepackage[utf8]{inputenc}
\usepackage{microtype}
\usepackage{inconsolata}
\usepackage{graphicx}
\usepackage{booktabs}

% Items that still need the author. Search the PDF for these before submitting.
\newcommand{\citeme}{\textbf{[CITE]}}
\newcommand{\unverified}{\textbf{[UNVERIFIED]}}
\newcommand{\snip}[1]{{\footnotesize\texttt{#1}}}

\title{Reproducing and Extending Zero-Shot Relation Extraction with REPaL}

\author{Manohar Naidu Bheesetti \\
  Student ID: 48975257 \\
  COMP8240: Applications of Data Science, Session 2, 2026 \\
  Macquarie University \\
  \texttt{manoharnaidu.bheesetti@students.mq.edu.au}}

\begin{document}
\maketitle

\begin{abstract}
REPaL trains a relation extractor from only a natural-language relation definition and an unlabelled corpus. I replicated it on the authors' DefOn-FewRel and DefOn-WikiZSL splits using their cached GPT-4o synthesis. Followup F1 was 69.62 on FewRel (paper: 74.61) and 45.95 on WikiZSL (paper: 47.80). I then applied it to SemEval-2010 Task 8 and NYT with locally served Qwen models, reaching 30.99 and 47.10 respectively in single-split runs. I also collected a new arXiv corpus with a six-relation schema and blank annotation sheets. The arXiv annotation, inter-annotator agreement, the REPaL run on arXiv, and the planned LLM-as-a-judge evaluation were not completed, and no results for them are reported.
\end{abstract}

\section{Paper Description and Justification}
\label{sec:paper}

Relation extraction (RE) identifies the semantic relation between two entities in text. Supervised RE needs labelled data for every relation of interest, and most zero-shot methods still need labelled data for other, ``seen'' relations. \citet{zhou2024grasping} study a harder setting, \emph{definition-only zero-shot RE}: the only inputs are a natural-language definition of the target relation and an unlabelled corpus. Negative relations are assumed unknown.

Their framework, REPaL, has three stages. (1) \emph{Definition-based seed construction}: an LLM writes positive examples from the definition, and initial negatives are sampled from the unlabelled corpus. (2) \emph{Pattern learning}: a small language model (SLM), here the NLI model \texttt{roberta-large-mnli}, derived from RoBERTa \citep{liu2019roberta}, is fine-tuned as a binary classifier in which the definition is the hypothesis and the sentence is the premise. (3) \emph{Feedback}: the SLM's predictions on the unlabelled corpus, together with the synthesis history, are returned to the LLM, which writes follow-up positive examples and negative relation definitions targeted at the SLM's errors. The SLM is then retrained.

Evaluation resembles cross-validation. With $R$ test relations there are $R$ binary tasks, each treating one relation as positive and the rest as negative. Precision, recall and F1 are computed per task and averaged. Because F1 is averaged per relation, it can be lower than both mean precision and mean recall.

The paper reports REPaL (GPT-4o) at P/R/F1 of 78.86/77.28/74.61 on DefOn-FewRel and 68.98/47.63/47.80 on DefOn-Wiki-ZSL, ahead of the zero-shot RoBERTa NLI baseline (F1 55.07 and 35.28). Its ablation shows that the LLM's initial positives contribute most: removing them gives F1 55.07 on DefOn-FewRel, while removing positive feedback gives 70.03 and removing negative feedback gives 73.04.

\paragraph{Justification.} The work was published at EMNLP 2024 (pp.\ 13462--13486) by authors from the University of Illinois Urbana-Champaign and the University of Virginia. EMNLP is ranked A* in the ICORE2026 ranking,\footnote{\url{https://portal.core.edu.au/conf-ranks/}, checked 2026-11-09.} and the paper had 30 citations on Semantic Scholar when checked on 2026-11-09. Feasibility was strong: code is public, and the authors released the evaluation splits and a cache of GPT-4o synthesis, so replication needs no API spending. Scientifically, a method that needs only a definition should in principle transfer to new relation schemas, which is what this project's extension tests.

\section{Original Datasets}
\label{sec:original}

Both benchmarks are derived by \citet{zhou2024grasping}. \textbf{DefOn-FewRel} comes from FewRel \citep{han2018fewrel}, whose public portion has 80 relations with 700 crowd-annotated instances each. \textbf{DefOn-WikiZSL} comes from Wiki-ZSL \citep{chen2021zs}, with 113 relations and 93,383 distantly supervised instances. Relations are Wikidata properties; the authors rewrote property definitions into sentences with entity placeholders. They sampled 5 groups of 14 FewRel relations and 3 groups of 15 Wiki-ZSL relations as test sets.

The paper states that the unlabelled corpus for each group is 10,000 instances down-sampled from the 899,996-instance corpus of \citet{gao2020neural}. The released splits differ, as Table~\ref{tab:splits} shows: every released \texttt{distant.json} has 100,000 instances. Section~\ref{sec:replication} returns to this.

\begin{table}[t]
\centering\small
\begin{tabular}{lrrr}
\toprule
Dataset & Rel. & Test (per split) & Pool \\
\midrule
FewRel (5 splits) & 14 & 9,796--9,799 & 100,000 \\
WikiZSL (3 splits) & 15 & 12,122--12,857 & 100,000 \\
\bottomrule
\end{tabular}
\caption{Released splits, counted from the JSON files.}
\label{tab:splits}
\end{table}

Each instance holds tokens plus head and tail mentions with token indices. Each relation has a definition prompt. Real examples from the test files follow.

\begin{quote}\small
\textbf{FewRel, P106 (occupation).} Definition: \snip{<ENT1> was/is the occupation of <ENT0> (a person)}.
Sentence: ``David Malo or Davida Malo ( 1793--1853 ) was a leading Native Hawaiian historian of the Kingdom of Hawaii .''
Head: \emph{david malo}; tail: \emph{historian}.

\textbf{WikiZSL, P101 (field of work).} Definition: \snip{<ENT1> (a field of work, academic area, or specialism) was/is the specialization of <ENT0> (a person or organization)}.
Sentence: ``Otto Toeplitz ( 1 August 1881 15 February 1940 ) was a German Jewish mathematician working in functional analysis .''
Head: \emph{Otto Toeplitz}; tail: \emph{functional analysis}.
\end{quote}

The authors attribute WikiZSL's lower scores to noise from distant supervision. These datasets are suitable because their definitions are precise and they have enough relations to average over, though the FewRel test instances are drawn from a single domain (Wikipedia).

\section{Replication}
\label{sec:replication}

\paragraph{Method.} I forked the authors' code\footnote{\url{https://github.com/ManoharNaidu/REPaL}} and replaced its two shell scripts, which contained unfilled placeholders and assumed four GPUs, with a Python driver that runs each dataset, split and stage, captures logs, and parses metrics. The hyperparameters follow the released configuration: \texttt{roberta-large-mnli}, learning rate $3\times10^{-5}$, 12 epochs, training batch 16, seed 3, and 15 positive plus 15 negative examples at both the initial and follow-up steps. Inference batch size was 512 instead of the authors' 3600--4200. Two stages were run: \emph{initial} (seed synthesis, training) and \emph{followup} (feedback-driven synthesis, retraining). The paper's headline number corresponds to \emph{followup}.

All eight splits used the authors' cached GPT-4o synthesis (\texttt{gpt-4o-2024-05-13}), so no new API calls were needed.

\paragraph{Environment and repairs.} Run logs show Python 3.12 in a Linux container, with runs started between 2026-10-06 and 2026-10-09. The logs do not record the exact torch, CUDA or GPU details. Notes taken on 2026-10-10 record Python 3.12.14, torch 2.11.0+cu128 and two RTX 5090 cards for the instance at that date, but package versions at run time were not recorded, so I treat these as indicative only. Several changes were needed for the 2024 code to run on current hardware: moving \texttt{BCELoss} outside \texttt{autocast}, passing \texttt{weights\_only=False} to \texttt{torch.load}, removing the torch pin (RTX 50-series cards need a newer CUDA build), and constraining \texttt{numpy<2} and \texttt{httpx<0.28}. I also capped a negative-feedback sampling step that crashed when a relation had too few confident predictions.

\paragraph{Results.} All 16 runs exited successfully. Tables~\ref{tab:replication} and~\ref{tab:splitf1} give the results. Each FewRel split took 160--197 minutes for the initial stage and about 97 minutes for followup. WikiZSL took 179--253 and 137--147 minutes.

\begin{table}[t]
\centering\small
\begin{tabular}{lrrr}
\toprule
 & P & R & F1 \\
\midrule
\multicolumn{4}{l}{\emph{DefOn-FewRel}} \\
Ours, initial & 67.04 & 75.51 & 64.88 \\
Paper, iter.\ 1 (Table 4) & 73.59 & 75.90 & 70.03 \\
Ours, followup & 79.42 & 68.81 & 69.62 \\
Paper, Table 1 & 78.86 & 77.28 & 74.61 \\
\midrule
\multicolumn{4}{l}{\emph{DefOn-WikiZSL}} \\
Ours, initial & 64.72 & 51.41 & 49.15 \\
Ours, followup & 68.69 & 42.33 & 45.95 \\
Paper, Table 1 & 68.98 & 47.63 & 47.80 \\
\bottomrule
\end{tabular}
\caption{Replication against the paper (means over 5 and 3 splits). The paper gives no initial-only WikiZSL figure.}
\label{tab:replication}
\end{table}

\begin{table}[t]
\centering\small
\begin{tabular}{lrr}
\toprule
Split & Initial F1 & Followup F1 \\
\midrule
FewRel 1 & 65.49 & 68.60 \\
FewRel 2 & 65.39 & 75.59 \\
FewRel 3 & 64.35 & 65.51 \\
FewRel 4 & 66.20 & 68.77 \\
FewRel 5 & 62.95 & 69.63 \\
\midrule
WikiZSL 1 & 51.25 & 51.27 \\
WikiZSL 2 & 45.91 & 39.26 \\
WikiZSL 3 & 50.29 & 47.32 \\
\bottomrule
\end{tabular}
\caption{Per-split F1. Sample standard deviation of followup F1 is 3.68 (FewRel) and 6.12 (WikiZSL).}
\label{tab:splitf1}
\end{table}

\paragraph{Assessment.} The replication is partial but close. WikiZSL followup F1 is 1.85 below the paper, within one standard deviation across splits. FewRel is 4.99 below, about 1.4 standard deviations. The pattern of the gap is informative. At followup, FewRel precision matches the paper (+0.56) while recall is 8.47 lower. At the initial stage the opposite holds: recall matches ($-$0.39) and precision is 6.55 lower. Feedback therefore moved our model further along the precision-recall trade-off than the paper reports. Feedback improved FewRel by 4.75 F1 points but reduced WikiZSL by 3.20, so the benefit reported in the paper was reproduced on only one of the two datasets.

\paragraph{Possible reasons for the gaps.} I tested none of these conclusively. (a) \emph{Pool size}: the paper specifies 10,000 unlabelled instances, whereas the release has 100,000. I built a 10,000-instance pool for FewRel splits 1 and 2 by seeded sampling from the released pool. Initial F1 rose from 65.49 to 73.58 and from 65.39 to 69.61, which suggests the pool matters. However, this is only two splits and one seed, and the split 1 run reused existing checkpoints for many relations rather than training from scratch, so it is not a clean measurement. (b) \emph{Software}: the code targeted Python 3.8; I used 3.12 and much newer libraries. (c) \emph{Variance}: I used one seed per split. (d) \emph{Pool composition}: my 10,000-instance sample comes from the released 100,000, not the authors' original corpus, which is not released.

\paragraph{Local LLM substitution.} OpenAI credits ran out during the extension work, so fresh synthesis moved to locally served models through an OpenAI-compatible vLLM server: Qwen2.5-32B-Instruct-AWQ \citep{qwen25} and, for one comparison, Qwen3-32B-AWQ \citep{qwen3}. The code needed no change. The 10,000-pool followup runs used Qwen2.5 on top of GPT-4o seeds (F1 65.90 and 70.27). Because these mix generators, I exclude them from the comparison with the paper. Results from different generators are never averaged together in this report.

\section{Construction of New Data}
\label{sec:newdata}

The proposal named SemEval-2010 Task 8 and BioCreative V CDR as existing datasets, and a newly built scientific-literature resource. BioCreative V CDR was not completed within the project timeline. NYT was added as a second existing dataset.

\subsection{SemEval-2010 Task 8}
SemEval-2010 Task 8 \citep{hendrickx2010semeval} labels the relation between two nominals. I obtained it through FewRel 2.0's \texttt{val\_semeval.json} \citep{gao2019fewrel}. The builder script removes the 1,864 ``Other'' instances and the 2 instances of Entity-Destination(e2,e1), leaving 8,851 instances and 17 directed relations (the proposal's table listed 19). Definitions are the task's guideline sentences, and a template makes each one direction-aware. Removing ``Other'' makes this easier in that respect than the official 19-way task.

\begin{quote}\small
\textbf{Cause-Effect(e1,e2).} Definition: \snip{<ENT0> (an event or object) yields <ENT1> (an effect)}.
Sentence: ``the singer , who performed three of the nominated songs , also caused a commotion on the red carpet .''
Head: \emph{singer}; tail: \emph{commotion}.
\end{quote}

\subsection{NYT}
NYT comes from FewRel 2.0's \texttt{val\_nyt.json}: 25 Wikidata relations with 100 news sentences each (2,500 instances). Of the 25 definitions, 18 reuse the REPaL authors' prompts and 7 were newly written in their style from Wikidata descriptions.

\begin{quote}\small
\textbf{P1056 (product or material produced).} Definition: \snip{<ENT1> (a product or material) was/is produced by <ENT0> (a business entity, government agency, industry, facility or process)}.
Head: \emph{Williams International}; tail: \emph{jet engines}, in a sentence about a ``demonstration by about 150 people against manufacturing of jet engines for cruise missiles at the Williams International Corporation plant''.
\end{quote}

\paragraph{Transductive construction.} For both datasets the builders write the \emph{same} test instances, with the relation grouping removed, as the unlabelled pool: 8,851 of 8,851 (SemEval) and 2,500 of 2,500 (NYT) pool instances are test instances. Labels are hidden, but pool-based mining and feedback see the test sentences. Test--pool overlap also exists in the released data, but is far smaller: 43.7\% in FewRel split 1, 5.9--6.7\% in splits 2--5, and under 1\% in WikiZSL. SemEval and NYT scores should therefore be read as transductive, and are probably optimistic relative to an unseen-sentence evaluation.

\subsection{New arXiv corpus}
This resource is intended to test REPaL on scientific text. Its \emph{collection and preprocessing are complete; its annotation is not}.

\paragraph{Collection.} A script queried the official arXiv API at no more than one request per 3.2 seconds. It sampled by month and primary category from 2025-01 to 2026-09, yielding 397 unique abstracts (cs.CL 83, cs.LG 80, cs.CV 80, cs.IR 80, stat.ML 74), dated 2025-01-01 to 2026-09-11.

\paragraph{Preprocessing.} Sentence splitting gave 2,947 sentences, of which 2,846 passed an 8--60 token filter (mean 28.4, median 27). A heuristic mention detector (acronyms, CamelCase names, metric words, scores, and phrases such as ``\ldots\ model'') found 1,384 sentences with two or more mentions, giving 3,737 candidate pairs in REPaL's \texttt{distant.json} format. No scientific NER model was used, so mentions are noisy.

\paragraph{Schema.} Table~\ref{tab:schema} shows the six relations, each stored in two directions (12 labels). The proposal's starting list changed: \texttt{proposes\_method} was dropped because its head, ``the paper'', is not an entity in the sentence, and \texttt{compares\_with\_baseline} became the directional \texttt{outperforms}.

\begin{table}[t]
\centering\footnotesize
\begin{tabular}{p{0.32\columnwidth}p{0.56\columnwidth}}
\toprule
Relation & Head $\rightarrow$ tail \\
\midrule
\texttt{evaluated\_on} & method/model $\rightarrow$ dataset/benchmark \\
\texttt{outperforms} & method $\rightarrow$ baseline/earlier method \\
\texttt{achieves\_result} & method $\rightarrow$ score/metric value \\
\texttt{builds\_on} & new method $\rightarrow$ earlier method \\
\texttt{applied\_to\_task} & method $\rightarrow$ task/problem/domain \\
\texttt{uses\_component} & method $\rightarrow$ technique/module/resource \\
\bottomrule
\end{tabular}
\caption{arXiv relation schema.}
\label{tab:schema}
\end{table}

\begin{quote}\small
\textbf{Unlabelled candidate pair.} From ``Large Language Models ( LLMs ) have demonstrated remarkable reasoning capabilities \ldots'': head \emph{LLMs}, tail \emph{discovering effective reasoning methods}. The tail shows the mention detector's noise.

\textbf{Definition.} \snip{<ENT0> (a method or model) is evaluated on <ENT1> (a dataset or benchmark)}.

\textbf{Annotation-sheet row ax002} (arXiv 2601.15239), label columns blank: ``Applied to contextualized word embeddings from language models, MCPCA maps stages of a debate on human nature, revealing a discussion between science and fiction over decades.''
\end{quote}

\paragraph{Annotation material and status.} Two identical sheets each hold 100 sentences: 20 per category, one sentence per abstract, sampled with keyword cues plus some cue-free sentences. Each sheet has blank columns for head entity, tail entity, relation and notes. An annotator guide, an agreement script based on Cohen's $\kappa$ \citep{cohen1960coefficient}, and a converter to REPaL's format exist, and have been tested only on made-up rows.

\textbf{No annotation took place.} All label cells in both sheets are empty. There are no annotators, no gold labels and no $\kappa$ value, and the arXiv test set was therefore never built.

\section{Results on New Data}
\label{sec:newresults}

Table~\ref{tab:newresults} gives the measured results. Each is a single split with seed 3, so none has a variance estimate.

\begin{table}[t]
\centering\small
\begin{tabular}{lrrr}
\toprule
 & P & R & F1 \\
\midrule
\multicolumn{4}{l}{\emph{SemEval, Qwen2.5-32B}} \\
Initial & 45.42 & 20.39 & 19.04 \\
Followup & 45.49 & 35.86 & 30.99 \\
\multicolumn{4}{l}{\emph{SemEval, Qwen3-32B}} \\
Initial & 42.55 & 30.97 & 25.10 \\
Followup & 46.72 & 25.70 & 19.64 \\
\multicolumn{4}{l}{\emph{SemEval, GPT-4o (archived run)}} \\
Initial & 57.18 & 34.59 & 35.67 \\
\multicolumn{4}{l}{\emph{NYT, Qwen2.5-32B}} \\
Initial & 51.95 & 56.96 & 48.35 \\
Followup & 68.65 & 45.68 & 47.10 \\
\bottomrule
\end{tabular}
\caption{Results on existing new datasets. The GPT-4o SemEval run has no followup stage because credits ran out.}
\label{tab:newresults}
\end{table}

\paragraph{SemEval.} SemEval is much weaker than the replication. Precision is moderate and recall is the main problem. An analysis of the existing logs and initial-stage inference files found that no relation model fires on 64.6\% of test instances. Direction confusion is not the main failure: 10.0\% of false positives come from the reverse-direction class. The two Component-Whole models produce 74\% of false positives, and Member-Collection(e1,e2) and Cause-Effect(e1,e2) have F1 near zero in all four runs analysed. The inference files come from a slightly different checkpoint from the reported numbers (recomputed macro P/R of 48.3/16.1 against reported 45.4/20.4), so these are patterns rather than exact figures.

A plausible cause is a mismatch between synthetic and real data. Real SemEval entities average 1.0--1.1 tokens in 19-token sentences. Seed entities (head/tail) average 1.9/2.1 tokens for GPT-4o in 23-token sentences, 3.6/3.7 for Qwen3 in 28-token sentences, and 4.2/4.9 for Qwen2.5 in 37-token sentences. Initial F1 follows the same order (35.67, 25.10, 19.04). This is a hypothesis from three single-seed runs, and within the Qwen2.5 run the seed span length does not track per-relation F1. Generator quality is also confounded with it.

\paragraph{NYT.} NYT reaches 47.10, close to WikiZSL (45.95), with a Qwen2.5 generator. Feedback shifted the balance from recall to precision and lowered F1 slightly.

\paragraph{Effect of feedback.} Followup improved on initial for FewRel (+4.75) and SemEval-Qwen2.5 (+11.95) but not WikiZSL ($-$3.20), NYT ($-$1.25), or SemEval-Qwen3 ($-$5.46). In the single-split cases, the differences may be within run-to-run variation, which I could not measure.

\paragraph{Relation to the replication.} REPaL's performance depends strongly on the dataset. Scores on the new data are not comparable with the replication as a clean test of generalisation, for three reasons: the generators differ (GPT-4o against Qwen), the SemEval and NYT evaluations are transductive, and each has one split and one seed. The results do show that definitions alone do not guarantee transfer: SemEval's short, abstract definitions and single-noun entities were hard for the pipeline, even with GPT-4o seeds.

\paragraph{Absent results.} There are \textbf{no results for the arXiv dataset}, because no gold labels exist and REPaL was not run on it. The planned \textbf{LLM-as-a-judge evaluation was not carried out}, so no judge scores or agreement figures exist. The BioCreative V CDR dataset was not completed within the project timeline. The per-iteration saturation analysis, a stretch goal, was also not done.

\section{Ethics}
\label{sec:ethics}

Following \citet{benotti-blackburn-2022-ethics}, I structure this section around research benefits, potential harms and vulnerable groups. Ethics statements in NLP papers have themselves become an object of study \citep{karamolegkou-etal-2025-ethical}, which supports treating them as substantive content. I cover the original paper first and then the issues that would arise if I published on my new data. The arXiv annotation was not completed, so the points about annotators below are prospective.

\paragraph{Research benefits.} REPaL needs only a relation definition and an unlabelled corpus, which could reduce the human annotation that RE otherwise requires. It also lowers the cost of building extractors for new relations. A replication adds value by reporting where results hold and where they do not. Here, that includes the mismatch between the paper's stated 10,000-instance pool and the released 100,000-instance pool.

\paragraph{Potential harms in the original paper.} The authors' ethics statement says that they do not verify the factualness of generated instances and ``simply count on the LLMs'', arguing that the text context is enough for sentence-level RE. They warn that adaptations needing factualness should be careful. I consider this a real limitation: seed examples in my logs name real historical figures, and those sentences are unverified. The paper also acknowledges that bias or errors can be introduced during synthesis. Its feedback stage targets the SLM's pattern biases, but that is not the same as addressing social bias. The evaluation data also contain sentences about named people, and FewRel relies on crowd workers' annotations.

\paragraph{Potential harms if I published on my new data.}
\emph{arXiv text reuse and author privacy.} Abstracts were collected through the official API within its rate limit. arXiv's terms permit use of descriptive metadata, defined to include titles and abstracts, under CC0 1.0 \citep{arxivtou}, but do not extend to storing and serving e-prints, and I collected no full text. Abstracts are not anonymous: the corpus keeps arXiv identifiers for attribution, and these link each sentence to its authors. Labels such as \texttt{outperforms}, attached to a named paper's work, could be read as evaluative claims about its authors. Text anonymisation methods such as differentially private rewriting \citep{meisenbacher-etal-2024-dp} exist, but would conflict with attribution and with the entity spans RE needs, so a release would instead need a takedown route for authors.

\emph{Annotator consent and risk.} Annotators would be participants in research, so the design would need assessment against the National Statement \citep{nhmrc-national-statement-2025}, including informed and voluntary consent, the freedom to withdraw, and proportionate risk. The risk appears low, as the task involves reading technical sentences, but the burden and any relationship between researcher and annotator would need attention. Because no annotators were recruited, there are no annotation data and no agreement figures, and I did not simulate any.

\emph{Responsible conduct and data management.} The Australian Code \citep{nhmrc-code-2018} frames responsible conduct and data management. In practice, I report shortfalls and non-clean runs, state the generator behind every result, and do not average across generators. Per-run configurations, logs and metrics are retained, and I keep a file manifest with checksums for the backup of results.

\emph{LLM-generated training data.} Synthetic examples can contain errors and encode the generator's biases, which the trained model can inherit. On SemEval, seed style differed clearly between GPT-4o and the Qwen models, showing that the generator can shape the data substantially. Local models also vary by version, so model identifiers are recorded for each run. Automated moral reasoning is an active research area \citep{pyatkin-etal-2023-clarifydelphi}, but I did not use such tools, and these judgements are my own.

\emph{Responsibility for AI-generated output.} Generative AI tools assisted with this project, as the unit permits. In line with Macquarie University's policy \citep{mq-ai-policy}, I remain responsible for this report and for checking AI-generated output.

\paragraph{Vulnerable groups.} Three groups deserve particular attention. People named in Wikipedia and news sentences cannot consent to automated extraction, and false relations could harm them. Authors, especially early-career researchers, could be affected by automated claims about their work. Prospective annotators could be vulnerable if recruited from students or others in a dependent relationship with the researcher, which is a further reason that consent would need to be genuinely voluntary.

\section{Conclusions and Reflections}
\label{sec:conclusions}

\paragraph{What replicated.} The released code, after repairs, ran end to end on all eight splits and two stages. Followup F1 was 69.62 on FewRel (paper: 74.61) and 45.95 on WikiZSL (paper: 47.80). The method's broad performance level replicated. FewRel's improvement from feedback also replicated, although with lower recall than reported. WikiZSL's did not: feedback lowered F1 by 3.20. The released unlabelled pools (100,000 instances) do not match the paper's stated 10,000, a reproducibility detail that a reader would miss without counting the files.

\paragraph{How well it generalised.} My primary research question asked whether REPaL generalises to new relation types and domains. The evidence is limited and mixed. NYT reached 47.10 and SemEval between 19.64 and 30.99 with local generators. An archived GPT-4o run reached 35.67 on SemEval at the initial stage. The scientific-literature domain, which was the most novel test, remains untested because the annotation was not completed. So I can say that REPaL's performance varied substantially across the two existing datasets, but I cannot claim that it generalises to scientific text. My second question, about which definitions give the best synthesis, was only addressed indirectly through the SemEval analysis.

\paragraph{What I would do differently.} First, I would schedule annotation much earlier and recruit annotators at the start, because it was the bottleneck that blocked the arXiv evaluation, the agreement study and part of the judge design. Second, I would hold the generator fixed across all comparisons and reserve credits for it, because the mid-project move from GPT-4o to local models confounded the extension results. Third, I would run several seeds on every new dataset, because single-split results cannot separate a real effect from noise. Fourth, I would hold out unseen sentences from the unlabelled pool on new datasets, avoiding a fully transductive evaluation. Finally, I would check dataset properties such as pool size against the paper before the main runs, not after them, and would log library versions and GPU details in every run, since I could not recover them afterwards.


\bibliographystyle{acl_natbib}
\bibliography{refs}

\end{document}
